\documentclass[reprint,amsmath,amssymb,aps,prapplied,longbibliography]{revtex4-2}

\usepackage{graphicx}
\usepackage{dcolumn}
\usepackage{bm}
\usepackage{physics}
\usepackage{hyperref}

\begin{document}

\title{Emergent Nematicity and Topological Defects in Discretized Spin Hamiltonians:\\ A BPS-Preserving Framework for Spatial Pattern Analysis}

\author{Author Name$^{1}$}
\email{author@institute.edu}
\affiliation{$^1$Department of Physics, University Name, City, Country}

\date{\today}

\begin{abstract}
Topological spin textures, such as domain walls and skyrmions, exhibit robust invariant properties governed by continuous field theories. However, simulating these continuous Hamiltonian systems on discrete square lattices inherently breaks the $SO(2)$ rotational symmetry, leading to severe lattice anisotropy. Here, we present a theoretical and computational framework that maps arbitrary spatially varying scalar fields (e.g., digitized $\mathbb{Z}_2$ images) onto a continuous $S^2$ spin manifold, serving as a highly inhomogeneous pinning field. We introduce a generalized spin Hamiltonian incorporating isotropic exchange and Dzyaloshinskii-Moriya interactions (DMI). By analytically coupling the finite-difference operators for the exchange and DMI, we preserve the Bogomolny-Prasad-Sommerfield (BPS) topological energy bound on the discrete lattice. Using a variational optimization approach anchoring the nearest-neighbor interactions, we perfectly restore the $SO(2)$ rotational invariance, reducing the grid-induced orientation error to $0.1^\circ \sim 0.2^\circ$. Through topological dyadic tensor projections of the emergent Néel-type chiral boundaries, our model autonomously generates macroscopic nematic director fields from the ground state configuration. We demonstrate that this physical framework provides an extremely robust methodology for detecting fractional ($q = \pm 1/2$) and integer ($q = \pm 1$) topological defects in complex structures such as geological strata and biometrics. Our results bridge the gap between topological condensed matter physics and algorithmic pattern recognition.
\end{abstract}

\maketitle

\section{Introduction}
In condensed matter physics, the interplay between continuous symmetries and discrete lattice geometries governs the formation of emergent topological structures. Chiral magnetic domain walls, skyrmions, and merons are stable topological excitations driven by the competition among exchange interactions, Dzyaloshinskii-Moriya interactions (DMI), and external magnetic fields. The topological protection of these textures ensures that their macroscopic invariants—such as the winding number and topological charge—remain highly robust against local structural fluctuations and thermal noise.

Recently, there has been a growing interest in translating these physical invariants into computational algorithms for spatial data analysis. Conventional algorithms for extracting geometric features from discrete $\mathbb{Z}_2$ scalar fields (e.g., digital images or pixelated sensors) suffer from severe aliasing and lattice anisotropy, as the point-group symmetry of the underlying grid (e.g., $C_4$ for a square lattice) discretizes continuous measurements. 

In this work, we demonstrate that mapping inhomogeneous spatial data onto an $S^2$ spin manifold allows for the exact extraction of morphological and topological information by utilizing the physics of magnetic relaxation. We treat the external spatial data as a spatially varying Zeeman pinning field. Under a tailored spin Hamiltonian, the sharp $\mathbb{Z}_2$ discontinuities relax into smooth, finite-width chiral domain walls. The core theoretical breakthrough of our work is the formulation of a discretized Hamiltonian where the exchange and DMI effective fields are explicitly coupled. This coupling strictly preserves the Bogomolny-Prasad-Sommerfield (BPS) energy bound and completely restores the continuous $SO(2)$ rotational symmetry on a square lattice when combined with a constrained variational optimization of the interaction constants.

Furthermore, we show that by projecting the in-plane components of the resulting Néel-type spin textures onto orthogonal anti-skyrmionic basis states, the model autonomously reconstructs macroscopic nematic director fields. This allows for the direct extraction of topological singularities ($q = \pm 1/2, \pm 1$) from highly noisy spatial fields, offering a paradigm shift from traditional computer vision towards physics-driven topological data analysis.

\section{Physical Model and Hamiltonian}

\subsection{Mapping to the $S^2$ Spin Manifold}
We consider a discrete 2D space where information is given by an inhomogeneous scalar field $H_z(x,y)$. To evaluate the geometry of this field, we introduce an auxiliary continuous 3D spin field $\mathbf{S}(x,y) = (S_x, S_y, S_z)$ subject to the constraint $|\mathbf{S}| = 1$. The external field $H_z(x,y)$ acts as an effective out-of-plane Zeeman interaction, imposing a boundary condition that maps the initial binary states (e.g., object and background) to the poles of the $S^2$ sphere ($S_z \to \pm 1$).

\subsection{Energy Minimization and Domain Wall Formation}
The system is governed by a classical Heisenberg-type continuous Hamiltonian incorporating DMI to induce chirality:
\begin{equation}
\mathcal{H} = \int d^2 r \left[ A (\nabla \mathbf{S})^2 + D \mathbf{S} \cdot (\nabla \times \mathbf{S}) - \mu_0 \mathbf{H}_{eff} \cdot \mathbf{S} \right],
\end{equation}
where $A$ is the exchange stiffness, $D$ is the DMI constant, and $\mathbf{H}_{eff} = H_z \hat{z}$. To achieve the ground state configuration, the spin system is relaxed via the steepest descent path, analogous to the overdamped limit of the Landau-Lifshitz-Gilbert (LLG) equation. 

During relaxation, the sharp boundary of $H_z(x,y)$ broadens into a domain wall of finite thickness $\delta \sim \sqrt{A/H_z}$. Because the energy is strictly localized within this domain wall, the contour integral of the boundary length transforms into an area integral of the localized energy density. The topological protection of the domain wall ensures that small local spatial fluctuations do not alter the macroscopic integrated energy, providing extraordinary robustness against noise.

\subsection{Topological Tensor Projections and Nematicity}
The presence of DMI enforces a Néel-type chiral structure at the boundaries, where the in-plane spin components $\mathbf{S}^\parallel = (S_x, S_y)$ align perpendicular to the domain wall. To extract the macroscopic orientation of the structure, we mathematically project this localized Néel state onto virtual topological basis states, specifically the $\alpha$-type and $\beta$-type anti-skyrmion fields, defined as $\mathbf{S}_{\alpha}^\parallel = (S_x, -S_y)$ and $\mathbf{S}_{\beta}^\parallel = (S_y, S_x)$.

The inner product of these fields generates a dyadic tensor field $\mathbf{T}$:
\begin{eqnarray}
T_x &=& \mathbf{S}^\parallel_{\text{Néel}} \cdot \mathbf{S}^\parallel_{\alpha} = S_x^2 - S_y^2, \\
T_y &=& \mathbf{S}^\parallel_{\text{Néel}} \cdot \mathbf{S}^\parallel_{\beta} = 2S_x S_y.
\end{eqnarray}
Because the tensor components satisfy $(S_x^2 - S_y^2, 2S_x S_y) \equiv (-\mathbf{S}) \otimes (-\mathbf{S})$, the resulting field is invariant under a $180^\circ$ rotation ($ \mathbf{S} \to -\mathbf{S} $). Consequently, the vector field $\mathbf{T}(x,y)$ manifests as a pure \textit{Nematic Director Field}. The global principal axis of the initial scalar field can thus be determined with high precision by integrating the tensor field: $\theta = \frac{1}{2} \arctan (\int T_y d^2r / \int T_x d^2r)$.

\section{Restoring Continuous Symmetries on the Lattice}

\subsection{BPS-Preserving Discretization}
Translating the continuous Hamiltonian to a discrete square lattice generally induces severe anisotropy due to the $C_4$ point group symmetry. To mitigate this, we define generalized interaction stencils (convolution kernels) $K_{D_x}$ and $K_{D_y}$ for the DMI gradient operators. 

Crucially, to maintain the Bogomolny-Prasad-Sommerfield (BPS) topological bound—where the energy can be minimized by completing the square of the first-order derivatives—the discrete exchange Laplacian operator must be analytically coupled to the DMI operators. We propose the element-wise coupling constraint:
\begin{equation}
K_{HJ} = K_{D_x}^{\circ 2} + K_{D_y}^{\circ 2}.
\label{eq:coupling}
\end{equation}
This coupling naturally restores the parity of the operators; the anti-symmetric DMI kernels (odd parity) are squared to form a positive, symmetric exchange kernel (even parity), ensuring that the range of both interactions remains perfectly matched. Furthermore, it structurally guarantees that the energy scaling respects $SO(2)$ invariance ($\cos^2\theta + \sin^2\theta = 1$) under grid rotations.

\subsection{Variational Optimization with Anchored Interactions}
While Eq.~(\ref{eq:coupling}) enforces the physical structure, the coefficients of the extended interaction neighborhood (e.g., next-nearest neighbors up to a $7\times7$ stencil) must be tuned to fully suppress the residual lattice anisotropy. We employ an L-BFGS-B variational optimization to minimize an objective function targeting the rotational variance of the geometric invariants.

To prevent the spin system from falling into non-physical singular states (the ``pixel-pinning" effect common in unregularized discrete spin models), we introduce an anchored regularization scheme. The nearest-neighbor interaction coefficients are fixed to their theoretically derived values, acting as an anchor, while higher-order interactions are variationally optimized. This physically constrained optimization yields a generalized Hamiltonian that acts completely isotropic on the square lattice, reducing the rotational orientation error to a negligible $0.1^\circ \sim 0.2^\circ$.

\section{Results: Topological Defects in Complex Textures}

To demonstrate the versatility of our physics-driven framework, we apply the model to complex, noisy grayscale textures, such as human fingerprints and geological layer formations. We conceptualize a grayscale texture as a 3D topographic manifold. By sectioning this manifold into discrete level-sets, we evaluate the spin relaxation independently on each layer. The superposition of the resulting dyadic tensor fields $\mathbf{T}_{th}$ leads to constructive interference along the principal morphological flow, while isotropic background noise undergoes destructive interference.

The final averaged field is a macroscopic nematic director field. In 2D nematic systems, points where the director field cannot be defined correspond to topological defects. We evaluate the 2D topological winding number around local domains:
\begin{equation}
q = \frac{1}{2\pi} \oint \nabla \phi \cdot dl,
\end{equation}
where $\phi$ is the angle of the nematic director. As observed in biological and liquid-crystal systems, our framework autonomously identifies half-integer defects ($q = \pm 1/2$, corresponding to minutiae bifurcations and layer branching) and integer defects ($q = \pm 1$, corresponding to cores and topological singularities). Unlike heuristic algorithmic feature extraction, our topological mapping is strictly invariant and extremely robust against pixel degradation. 

\section{Conclusion}
We have introduced a generalized framework utilizing continuous spin Hamiltonians on discrete grids to extract morphological and topological invariants from generic scalar fields. By analytically coupling the effective exchange and DMI kernels and utilizing an anchored variational optimization, we have overcome the fundamental lattice anisotropy of square grids, preserving the continuous BPS bound and $SO(2)$ symmetry. The autonomous emergence of nematic director fields and the exact identification of fractional and integer topological defects demonstrate the immense potential of applying condensed matter principles to information theory, computational geometry, and pattern recognition. Future studies may extend this framework to 3D voxel spaces to detect more complex topological invariants, such as Hopfions, in volumetric medical or material imaging.

\begin{acknowledgments}
The authors acknowledge insightful discussions with colleagues. This research was supported by [Grant Information].
\end{acknowledgments}

\end{document}